% Generated draft: requires mathematical and source review.
\documentclass[11pt]{article}
\usepackage[margin=25mm]{geometry}
\usepackage{fontspec}
\setmainfont{Latin Modern Roman}
\usepackage{amsmath,amssymb,mathtools,unicode-math}
\setmathfont{Latin Modern Math}
\usepackage{enumitem,needspace,xurl,hyperref}
\setlength{\parindent}{0pt}
\setlength{\parskip}{5pt}
\setlength{\emergencystretch}{3em}
\newcommand{\N}{\mathbb{N}}
\newcommand{\Z}{\mathbb{Z}}
\newcommand{\Q}{\mathbb{Q}}
\newcommand{\R}{\mathbb{R}}
\newcommand{\C}{\mathbb{C}}
\newcommand{\F}{\mathbb{F}}
\newcommand{\A}{\mathbb{A}}
\newcommand{\PP}{\mathbb{P}}
\newcommand{\E}{\mathbb{E}}
\newcommand{\eps}{\varepsilon}
\newcommand{\dd}{\,\mathrm{d}}
\begin{document}
\section*{UnsolvedMath ID: 30002850}
\textbf{Source problem code:} catalogue entry
\subsection{Definitions and mathematical statement}
Let \(G\) be a finite group and let \(P\in\operatorname{Syl}_2(G)\), so that \(P\) has order equal to the largest power of \(2\) dividing \(|G|\). Its normalizer is
\[
N_G(P)=\{g\in G:gPg^{-1}=P\}.
\]
The subgroup \(P\) is self-normalizing if \(N_G(P)=P\). Let \(\operatorname{Irr}(G)\) denote the set of irreducible complex characters of \(G\), and put
\[
\operatorname{Irr}_{2'}(G)=\{\chi\in\operatorname{Irr}(G):2\nmid\chi(1)\}.
\]
The McKay conjecture at the prime \(2\) is the assertion that
\[
\bigl|\operatorname{Irr}_{2'}(G)\bigr|
=
\bigl|\operatorname{Irr}_{2'}(N_G(P))\bigr|
\]
for every finite group \(G\) and every \(P\in\operatorname{Syl}_2(G)\). In the self-normalizing case, the assertion to be reduced is therefore
\[
\text{for every finite }G\text{ and every }P\in\operatorname{Syl}_2(G),\qquad
N_G(P)=P\ \Longrightarrow\
\bigl|\operatorname{Irr}_{2'}(G)\bigr|
=
\bigl|\operatorname{Irr}_{2'}(P)\bigr|.
\]
A finite group \(H\) is almost simple if there is a finite nonabelian simple group \(S\) such that
\[
S\leq H\leq\operatorname{Aut}(S),
\]
where \(S\) is identified with its group of inner automorphisms.

The supplied question asks whether the McKay conjecture for finite groups with self-normalizing Sylow \(2\)-subgroups can be directly reduced to almost simple groups. The displayed universally quantified assertion identifies the target of such a reduction; it does not specify the reduction hypotheses. In particular, the record does not say which almost simple groups are to be considered, whether their Sylow \(2\)-subgroups must also be self-normalizing, or what additional conditions constitute a direct reduction. The supplied dated assessment reports that the underlying sentence is incomplete. Consequently, a precise reduction theorem cannot be recovered from this record without further source clarification.
\subsection{Short English statement}
The question asks whether the McKay equality counting irreducible complex characters of odd degree, for every finite group with a self-normalizing Sylow \(2\)-subgroup, can be established by a direct reduction to almost simple groups. The intended reduction and its hypotheses are missing. The supplied open label is dated curation, and this statement-only draft requires human mathematical review.
\subsection{Sources}
\begin{itemize}
\item Local UnsolvedMath registry, permanent ID \(30002850\), source problem code \texttt{catalogue entry}, title: \emph{Direct McKay Reduction for Self-Normalizing Sylow Two-Subgroups}. The record belongs to the Oberwolfach Reports Open Problems collection and gives the proposed year as \(2015\); no proposer is named.
\item Supplied report citation: \emph{Representations of Finite Groups (2016)}, Oberwolfach Reports, pp. 971--1027, DOI \texttt{10.4171/owr/2015/18}. Supplied source URL and literature source labelled ``original OWR report'': \url{https://doi.org/10.4171/owr/2015/18}. The collection describes these reports as records of mathematical workshops and their open-problem discussions. This citation has not been independently consulted here.
\item Literature assessment dated \texttt{2026-08-21}: the supplied status is open, but the assessment says that the sentence ends before stating the proposed direct reduction, leaving the exact claim incomplete and unverifiable. This is supplied dated curation, not an independently verified literature assessment. The record explicitly calls for independent verification of time-sensitive status and cited claims.
\item The registry also labels the statement ``exact'' and describes it as clear, self-contained, publication-ready, and retained byte-for-byte. These supplied verification claims conflict with the supplied assessment of incompleteness and do not establish the missing reduction statement.
\end{itemize}
\end{document}
